\documentclass[11pt,a4paper]{article}

\input{../../sablony/pripravy/priprava.sty}

\nastavPripravu{4. ročník}{Kombinatorika}{35}{Variace II -- omezení a smíšené úlohy}

\begin{document}

\blok{Motivace a návaznost}

Když je první znak čísla nenulový nebo některé osoby nesmějí být zvoleny, nelze bez rozmyslu použít jeden hotový vzorec. Propojujeme postupné volby s permutacemi.

\poznamka{Předpoklady: permutace, variace, pravidla součtu a součinu. Cíl: rozpoznat kombinaci omezení, rozdělit případy a zvolit model v úlohách smíšených s permutacemi.}

\blok{Výklad a odvození}

U čísel nesmí být první cifra nula; u kódů smí, pokud zadání neříká jinak. U omezených výběrů může být snazší zafixovat určitou pozici, vytvořit blok, odečíst zakázané výsledky nebo rozdělit případy na nepřekrývající se části.
\begin{definiceBox}
\textbf{Doplňková metoda} počítá přípustné výsledky jako všechny výsledky minus zakázané výsledky. Všechny případy musíme započítat právě jednou.
\end{definiceBox}

\blok{Společně řešené příklady}

\begin{prikladReseniBox}[title={Příklad 1}]
Kolik čtyřciferných čísel s navzájem různými ciframi lze vytvořit z číslic $0,1,2,3,4,5$?
\tcblower
První cifra má $5$ možností; další postupně $5,4,3$. Celkem $5\cdot5\cdot4\cdot3=300$.
\end{prikladReseniBox}



\begin{prikladReseniBox}[title={Příklad 2}]
Ze sedmi různých studentů obsadíme předsedu, zástupce a zapisovatele. Zapisovatelem nesmí být Petr. Kolik možností?
\tcblower
Zapisovatel: šest možností. Předseda: šest ze zbylých, zástupce pět. Celkem $6\cdot6\cdot5=180$.
\end{prikladReseniBox}



\begin{prikladReseniBox}[title={Příklad 3}]
Na slavnost vystupuje sedm různých hudebních skupin, jedna musí být poslední. Kolik pořadí je možných? Jak se úloha liší od výběru prvních tří vystupujících?
\tcblower
Pevné poslední místo: $6!=720$ permutací. Při výběru tří různých prvních míst ze sedmi jde o $V(7,3)=210$ variací.
\end{prikladReseniBox}



\blok{Samostatná práce}

\begin{zadaniReseniBox}[title={Samostatná práce 1}]
Kolik trojciferných čísel bez opakování vznikne z číslic $0,1,2,3,4$?
\tcblower
$4\cdot4\cdot3=48$.
\end{zadaniReseniBox}



\begin{zadaniReseniBox}[title={Samostatná práce 2}]
Z osmi osob vybereme kapitána a zástupce. Kapitánem musí být Anna nebo Boris. Kolik obsazení?
\tcblower
Kapitán má dvě možnosti, zástupce sedm: $2\cdot7=14$.
\end{zadaniReseniBox}



\begin{zadaniReseniBox}[title={Samostatná práce 3}]
Seřaď šest různých knih tak, aby dvě konkrétní stály na prvních dvou místech v libovolném pořadí.
\tcblower
$2!\cdot4!=48$.
\end{zadaniReseniBox}



\blok{Rozšíření}

\begin{zadaniReseniBox}[title={Rozšíření 1}]
Kolik čtyřciferných sudých čísel bez opakování lze sestavit ze všech deseti číslic $0,\ldots,9$?
\tcblower
Poslední cifra 0: $9\cdot8\cdot7=504$. Poslední cifra 2,4,6,8: $4\cdot8\cdot8\cdot7=1792$. Celkem $2296$.
\end{zadaniReseniBox}



\blok{Domácí úkol}

\poznamka{Doporučená práce doma: přibližně 15--30 minut. Úlohy jsou seřazeny od základních k náročnějším; učitel může poslední podúlohu označit jako volitelnou.}

\begin{zadaniReseniBox}[title={Domácí úkol}]

\begin{enumerate}[label=\alph*),leftmargin=*,itemsep=2pt]

\item Kolik čtyřmístných kódů bez opakování lze vytvořit ze sedmi různých symbolů?
\item Kolik čtyřciferných čísel bez opakování vytvoříme z číslic $0,1,2,3,4,5$?
\item Ze šesti osob obsadíme tři funkce, přičemž předsedou smí být jen Anna nebo Eva. Kolik možností?
\item Ve frontě stojí šest různých osob; Petra musí stát první a Adam poslední. Urči počet pořadí a vysvětli, proč jde o permutaci, ne variaci.
\item Z číslic $1,2,3,4,5,6$ vytvoř čtyřciferná lichá čísla bez opakování. Kolik jich je?

\end{enumerate}

\tcblower

\begin{enumerate}[label=\alph*),leftmargin=*,itemsep=2pt]

\item $7\cdot6\cdot5\cdot4=840$.
\item $5\cdot5\cdot4\cdot3=300$.
\item $2\cdot5\cdot4=40$.
\item $4!=24$; všech šest osob se řadí.
\item Jednotky 1,3,5: $3$ možnosti; zbylá místa $5\cdot4\cdot3$, celkem $180$.

\end{enumerate}

\end{zadaniReseniBox}

\poznamka{Orientační práce v hodině: 0--10 min opakování; 10--28 min rozklad podle omezení; 28--40 min smíšené úlohy; 40--45 min kontrola.}

\end{document}
